% This file provides examples of some useful macros for typesetting
% dissertations.  None of the macros defined here are necessary beyond
% for the template documentation, so feel free to change, remove, and add
% your own definitions.
%
% We recommend that you define macros to separate the semantics
% of the things you write from how they are presented.  For example,
% you'll see definitions below for a macro \file{}: by using
% \file{} consistently in the text, we can change how filenames
% are typeset simply by changing the definition of \file{} in
% this file.
% 
%% The following is a directive for TeXShop to indicate the main file
%%!TEX root = diss.tex

\newcommand{\NA}{\textsc{n/a}}	% for "not applicable"
\newcommand{\eg}{e.g.,\ }	% proper form of examples (\eg a, b, c)
\newcommand{\ie}{i.e.,\ }	% proper form for that is (\ie a, b, c)
\newcommand{\etal}{\emph{et al}}

% Some useful macros for typesetting terms.
\newcommand{\file}[1]{\texttt{#1}}
\newcommand{\class}[1]{\texttt{#1}}
\newcommand{\latexpackage}[1]{\href{http://www.ctan.org/macros/latex/contrib/#1}{\texttt{#1}}}
\newcommand{\latexmiscpackage}[1]{\href{http://www.ctan.org/macros/latex/contrib/misc/#1.sty}{\texttt{#1}}}
\newcommand{\env}[1]{\texttt{#1}}
\newcommand{\BibTeX}{Bib\TeX}

% Define a command \doi{} to typeset a digital object identifier (DOI).
% Note: if the following definition raise an error, then you likely
% have an ancient version of url.sty.  Either find a more recent version
% (3.1 or later work fine) and simply copy it into this directory,  or
% comment out the following two lines and uncomment the third.
\DeclareUrlCommand\DOI{}
\newcommand{\doi}[1]{\href{http://dx.doi.org/#1}{\DOI{doi:#1}}}
%\newcommand{\doi}[1]{\href{http://dx.doi.org/#1}{doi:#1}}

% Useful macro to reference an online document with a hyperlink
% as well with the URL explicitly listed in a footnote
% #1: the URL
% #2: the anchoring text
\newcommand{\webref}[2]{\href{#1}{#2}\footnote{\url{#1}}}

% epigraph is a nice environment for typesetting quotations
\makeatletter
\newenvironment{epigraph}{%
	\begin{flushright}
	\begin{minipage}{\columnwidth-0.75in}
	\begin{flushright}
	\@ifundefined{singlespacing}{}{\singlespacing}%
    }{
	\end{flushright}
	\end{minipage}
	\end{flushright}}
\makeatother
% Macros for noting things needing to be changed.
% The following definitions will also output a warning to the console

\newcommand{ \fixme}[1]{\typeout{**FIXME** #1}\textbf{[fix: #1]}}
\newcommand{\awk}{\typeout{**AWKWARD**}\textbf{[awk]}}
\newcommand{\citeme}{\typeout{**CITATION**}\textbf{[cite]}}
\newcommand{\refme}{\typeout{**REFER**}\textbf{[ref]}}
\newcommand{\bonus}[1]{\typeout{**BONUS**}\textbf{[Bonus: #1]}}

%Theorems, definitions and other stuff. This ares is finicky, don't touch!

\newtheorem{thm}{Theorem}
\numberwithin{thm}{chapter}

\newaliascnt{defn}{thm}
\newtheorem{defn}[defn]{Definition}
\aliascntresetthe{defn}
\providecommand*{\defnautorefname}{Definition}

\newaliascnt{lem}{thm} 
\newtheorem{lem}[lem]{Lemma}
\aliascntresetthe{lem}
\providecommand*{\lemautorefname}{Lemma}

\newaliascnt{cor}{thm} 
\newtheorem{cor}[cor]{Corollary}
\aliascntresetthe{cor}
\providecommand*{\corautorefname}{Corollary}

\newaliascnt{prop}{thm} 
\newtheorem{prop}[prop]{Proposition}
\aliascntresetthe{prop}
\providecommand*{\propautorefname}{Proposition}

\theoremstyle{definition}
\newaliascnt{exa}{thm}
\newtheorem{exa}[exa]{Example}
\aliascntresetthe{exa}
\providecommand*{\exaautorefname}{Example}

% Useful commands for shorter writing

\newcommand{\en}{\mathbb{N}}
\newcommand{\ze}{\mathbb{Z}}
\newcommand{\re}{\mathbb{R}}
\newcommand{\co}{\mathbb{C}}
\newcommand{\qu}{\mathbb{Q}}
\newcommand{\oo}{\mathbb{0}}
\newcommand{\one}{\mathbb{1}}

\newcommand{\del}{\partial}

\newcommand{\mc}[1]{\mathcal{#1}}
\newcommand{\mon}[2]{#1_1\cdots#1_#2}
\newcommand{\row}[2]{#1_1,\ldots, #1_#2}
\newcommand{\fullx}[3]{#1_{[#2]}^#3}
\newcommand{\fullox}[3]{\ol{#1}_{[#2]}^#3}

\newcommand{\pring}[3]{#1[\row #2#3]}
\newcommand{\lring}[3]{#1[#2_1,\ol{#2_1},\ldots,#2_#3,\ol{#2}_#3]}
\newcommand{\psring}[3]{#1[[\row #2#3]]}
\newcommand{\lsring}[3]{#1((\row #2#3))}

\newcommand{\Rpring}[2]{R[#1_1,\ldots,#1_#2]}
\newcommand{\Rlring}[2]{R[#1_1,\ol{#1_1},\ldots,#1_#2,\ol{#1}_#2]}
\newcommand{\Rpsring}[2]{R[[#1_1,\ldots,#1_#2]]}
\newcommand{\Rlsring}[2]{R((#1_1,\ldots,#1_#2))}


\newcommand{\deff}{\mathop{\mathrm{def}}}
\newcommand{\mdeff}{\mathop{\mathrm{mdef}}}
\newcommand{\supp}{\mathop{\mathrm{supp}}}
\newcommand{\orb}{\mathrm{Orbit}}
\newcommand{\stab}{\mathop{\mathrm{Stab}}}

\newcommand{\ol}[1]{\overline{#1}}

\newcommand{\qbinom}[2]{\begin{bmatrix}#1\\#2\end{bmatrix}}

\newcommand{\um}{\scalebox{0.5}[1.0]{\( - \)}}
\newcommand{\Q}[2]{Q_{#1}^{#2}}
\newcommand{\Qf}[1]{Q_{1,2}^{#1}}
\newcommand{\Qc}[1]{\wc{Q}_{1,2}^{#1}}
\newcommand{\Qh}[1]{\wh{Q}_{1,2}^{#1}}
\newcommand{\Qch}[1]{\wc{\wh{Q}}_{1,2}^{#1}}
\newcommand{\es}[2]{S_{#1}^{#2}}
\newcommand{\Sf}[1]{S_{1,2}^{#1}}
\newcommand{\aye}[2]{A_{#1}^{#2}}
\newcommand{\ext}[1]{#1[#1]}
\newcommand{\sign}{\mathop{\mathrm{sign}}}

\newcommand{\wh}{\widehat}
\newcommand{\wc}{\widecheck}
\newcommand{\wch}[1]{\wc{\wh #1}}
\newcommand{\wcal}[1]{\widetilde{\mc{#1}}}

\newcommand{\com}{\,,\quad}
\newcommand{\qandq}{\quad\text{and}\quad}
\newcommand{\qorq}{\quad\text{or}\quad}

\newcommand{\exx}{\!\!}
\newcommand{\exxx}{\!\!\!}
\newcommand{\exxxx}{\!\!\!\!}
\newcommand{\exxxxx}{\!\!\!\!\!}

\newcommand{\classic}{classical }
\newcommand{\Classic}{Classical }
\newcommand{\ug}{U}
\newcommand{\oneD}{one-dimensional }
\newcommand{\twoD}{two-dimensional }
\newcommand{\highD}{higher-dimensional }
\newcommand{\lowD}{lower-dimensional }
\newcommand{\dD}{$d$-dimensional }
\newcommand{\OneD}{One-dimensional }
\newcommand{\TwoD}{Two-dimensional }
\newcommand{\HighD}{Higher-dimensional }
\newcommand{\highSym}{highly-symmetric }
\newcommand{\HighSym}{Highly-symmetric }
\newcommand{\akm}{algebraic kernel method }
\newcommand{\akms}{algebraic kernel methods }
\newcommand{\okm}{obstinate kernel method }
\newcommand{\okms}{obstinate kernel methods }

\newcommand{\textitbf}[1]{\textit{\textbf{#1}}}


% DRAWING STUFF

\newcommand{\modelgraph}[8]{
\raisebox{-0.35\height}{\begin{tikzpicture}[scale=0.4]
	\draw [white] (-1.15,-1.15) rectangle (1.15,1.15);
	\draw (0,0) -- (-#1,#1);
	\draw (0,0) -- (0,#2);
	\draw (0,0) -- (#3,#3);
	
	\draw (0,0) -- (-#4,0);
	\draw (0,0) -- (#5,0);
	
	\draw (0,0) -- (-#6,-#6);
	\draw (0,0) -- (0,-#7);
	\draw (0,0) -- (#8,-#8);
\end{tikzpicture}
}}

\newcommand{\modelgraphB}[8]{
\raisebox{-0.5\height}{\begin{tikzpicture}[scale=0.7]
	\draw [white] (-1.15,-1.15) rectangle (1.15,1.15);
	\draw (0,0) -- (-#1,#1);
	\draw (0,0) -- (0,#2);
	\draw (0,0) -- (#3,#3);
	
	\draw (0,0) -- (-#4,0);
	\draw (0,0) -- (#5,0);
	
	\draw (0,0) -- (-#6,-#6);
	\draw (0,0) -- (0,-#7);
	\draw (0,0) -- (#8,-#8);
\end{tikzpicture}
}}

\newcommand{\modelgraphC}[3]{
\raisebox{-0.5\height}{\begin{tikzpicture}[scale=0.7]
	\draw [white] (-1.15,0) rectangle (1.15,1.15);

	\draw (0,0) -- (0,#1);
	\draw (0,0) -- (#2,#2);
	\draw (0,0) -- (#3,0);
	
\end{tikzpicture}
}}

\newcommand{\cartmodel}{\modelgraph 01011010}
\newcommand{\diagmodel}{\modelgraph10100101}
\newcommand{\fullmodel}{\modelgraph 11111111}

% END
